\documentclass[12pt,a4paper]{amsart}

\usepackage{amsmath,amssymb,amsthm}
\usepackage{mathtools}
\usepackage{enumitem}
\usepackage{hyperref}

\newtheorem{theorem}{Theorem}[section]
\newtheorem{lemma}[theorem]{Lemma}
\newtheorem{proposition}[theorem]{Proposition}
\newtheorem{corollary}[theorem]{Corollary}
\theoremstyle{definition}
\newtheorem{definition}[theorem]{Definition}
\newtheorem{problem}[theorem]{Open Problem}
\newtheorem{conjecture}[theorem]{Conjecture}
\theoremstyle{remark}
\newtheorem{remark}[theorem]{Remark}

\newcommand{\Rig}{\textup{[Rig]}}
\newcommand{\Num}{\textup{[Num]}}
\newcommand{\Warn}{\textup{[Warning]}}
\newcommand{\Open}{\textup{[Open]}}

\begin{document}

\title{Mini-Model for the Fold Amplitude:\\
  Phase Modulation vs.\ Singularity Order,\\
  and the Sharpness of Ces\`aro}

\author{[Author]}

\begin{abstract}
This note studies
\[
  \Phi(u):=\int_{\mathbb{R}}
  |\beta|^{-1/2}\chi(\beta)
  e^{i(\beta u + c_3\beta^3)}\,d\beta,
  \qquad c_3>0,
\]
and $F(U):=\int_0^U\Phi(u)\,du$.
The following is established:
\begin{enumerate}[\rm(1)]
\item \textbf{Rigorous upper bound:}
  $|\Phi(u)|=O(|u|^{-1/2})$
  (Riemann--Lebesgue).
\item \textbf{Numerical evidence:}
  fitted exponent
  $|\Phi(u)|\approx|u|^{-0.53}$,
  $|F(U)|\approx U^{0.48}$;
  consistent with $|u|^{-1/2}$,
  no new scaling regime observed.
\item \textbf{Correct interpretation:}
  the cubic term acts as a
  high-frequency phase modulation;
  it does not shift the
  Fourier-singularity order
  of $|\Phi(u)|$.
\item \textbf{Open conjecture:}
  $\limsup_{|u|\to\infty}
  |u|^{1/2}|\Phi(u)|>0$
  (Ces\`aro bound is sharp).
  \emph{Not proved;
  numerically supported.}
\end{enumerate}
The $S+P$ decomposition
(singularity + saddle)
is used heuristically;
the two mechanisms are
\emph{not} asymptotically
orthogonal and cannot be
treated as independent
summands without a
formal partition of unity.
\end{abstract}

\maketitle

\section{What Is Rigorous}
\label{sec:rigorous}

\begin{lemma}[Upper bound \Rig]
\label{lem:upper}
$|\Phi(u)|=O(|u|^{-1/2})$
and $|F(U)|=O(U^{1/2})$.
Proof: $|\beta|^{-1/2}\chi\in L^1$;
Riemann--Lebesgue applies.
The cubic factor
$e^{ic_3\beta^3}$ does not
change $L^1$ membership;
the bound is independent
of $c_3$.
\end{lemma}

\begin{remark}[What the bound
  does not say \Rig]
\label{rem:not-say}
The upper bound $O(|u|^{-1/2})$
does not imply:
\begin{itemize}[nosep]
\item $|\Phi(u)|=\Theta(|u|^{-1/2})$
  (lower bound is open);
\item that Ces\`aro is sharp;
\item that the cubic term
  is irrelevant.
\end{itemize}
\end{remark}

\section{The Two Mechanisms
  and Their Limitation}
\label{sec:mechanisms}

\begin{remark}[Formal S+P split
  \Warn]
\label{rem:SP}
Two mechanisms formally
contribute to $\Phi(u)$:
\begin{enumerate}[\rm(S)]
\item[(S)] \textbf{Singularity:}
  the endpoint $\beta=0$
  contributes
  $\sim C_0 |u|^{-1/2}
  e^{i\pi/4}$
  (Fresnel; $C_0\ne 0$).
\item[(P)] \textbf{Saddle}
  (for $u<0$):
  $\beta_0=\sqrt{|u|/(3c_3)}$
  contributes formally
  $\sim C'|u|^{-1/2}
  e^{i\phi(\beta_0)}$
  with $\phi(\beta_0)
  =-\frac{2}{3}(|u|/3c_3)^{3/2}$.
\end{enumerate}
\textbf{Critical warning:}
S and P are \emph{not}
orthogonal asymptotically.
The decomposition
$\Phi=\Phi_S+\Phi_P$
is not canonical:
both contributions arise
from the same integral
over the same $\beta$-axis;
a split requires a
partition of unity
$\chi_S+\chi_P=1$
(separating a neighbourhood
of $\beta=0$ from one of
$\beta_0$) with explicit
error bounds in $u$.
Without this,
``$\Phi_S+\Phi_P$''
is a formal decomposition,
not an asymptotic identity.
\end{remark}

\begin{remark}[Saddle scale
  and separation \Rig]
\label{rem:scale}
The saddle
$\beta_0=\sqrt{|u|/(3c_3)}
\to\infty$ as $|u|\to\infty$.
For $|u|\gg 1$, the saddle
is far from $\beta=0$;
a partition
$\chi_S(\beta):=
\chi(\beta/\beta_0^{1/2})$,
$\chi_P:=1-\chi_S$
would separate the two
regions with explicit
error $O(\beta_0^{-N})$.
This partition would
make the S+P split rigorous,
but the computation of the
error in $u$ has not been
carried out here.
\end{remark}

\section{The Correct Conclusion:
  Phase Modulation}
\label{sec:phase}

\begin{proposition}[Phase
  modulation, not order shift
  \Rig+\Num]
\label{prop:phase}
The effect of the cubic
term $e^{ic_3\beta^3}$
in~\eqref{eq:Phi} is:
\begin{enumerate}[\rm(i)]
\item \textbf{On the phase:}
  the saddle phase
  $\phi(\beta_0)
  =-\frac{2}{3}(|u|/3c_3)^{3/2}$
  introduces a rapidly
  oscillating factor
  $e^{-i\frac{2}{3}(|u|/3c_3)^{3/2}}$
  in the saddle contribution;
  this has frequency
  $\sim|u|^{1/2}\to\infty$.
\item \textbf{On the amplitude:}
  numerically, no shift
  in the scaling exponent
  from $-1/2$ is observed;
  the upper bound $O(|u|^{-1/2})$
  is not improved.
\item \textbf{Correct interpretation:}
  the cubic term acts as
  a \emph{high-frequency
  phase modulation};
  it does not change
  the Fourier-singularity
  order of $\Phi$.
\end{enumerate}
\end{proposition}

\begin{remark}[Why phase alone
  cannot improve amplitude
  \Rig]
\label{rem:phase-alone}
If $|\Phi(u)|=\Theta(|u|^{-1/2})$
(i.e., the lower bound holds),
then $\Phi(u)=C_*(u)
|u|^{-1/2}$ where
$C_*(u)$ oscillates
but satisfies
$\limsup|C_*(u)|>0$.
The rapid oscillation of
$C_*(u)$ means that on
average $C_*(u)$ does not
systematically cancel;
but it also means that
on a sparse sequence
$|C_*(u_n)|$ might be small.
The question
$\limsup|C_*(u)|>0$
cannot be answered
from the upper bound alone;
it requires the lower
bound analysis.
\end{remark}

\section{Numerical Evidence
  and Its Limits}
\label{sec:numerical}

\begin{proposition}[Numerical
  findings \Num]
\label{prop:num}
With $c_3=1$, Gaussian
cutoff, $\beta\in[-50,50]$:
\[
  |\Phi(u)|\approx|u|^{-0.53},
  \quad
  |F(U)|\approx U^{0.48}.
\]
The deviation from
$-0.5$ and $0.5$ is
within expected numerical
error (finite integration
range, discrete sampling).
No scaling regime
consistent with $u^{-1/3}$,
$u^{-3/4}$, or $U^{2/3}$
is observed.
\end{proposition}

\begin{remark}[Limits of
  the numerical evidence
  \Warn]
\label{rem:num-limits}
The numerical computation:
\begin{itemize}[nosep]
\item supports $O(|u|^{-1/2})$
  as the correct order;
\item does \emph{not} prove
  $\Omega(|u|^{-1/2})$
  (a lower bound requires
  analytic argument,
  not sampling);
\item uses a specific
  Gaussian cutoff;
  the actual PSWF amplitude
  may differ;
\item is restricted to
  $|u|\le 30$;
  the asymptotic regime
  $|u|\to\infty$ is
  extrapolated, not observed.
\end{itemize}
\end{remark}

\section{The Single Open Problem}
\label{sec:open}

\begin{conjecture}[Ces\`aro is
  sharp \Open]
\label{conj:sharp}
\[
  \limsup_{|u|\to\infty}
  |u|^{1/2}|\Phi(u)|>0.
\]
Equivalently:
$|F(U)|=\Theta(U^{1/2})$
(Ces\`aro is not improvable
for this model).
\end{conjecture}

\begin{remark}[Why this is
  hard \Open]
\label{rem:hard}
The conjecture requires
showing that the rapidly
oscillating coefficient
$C_*(u)=|u|^{1/2}\Phi(u)$
does not decay to zero
along all sequences.
A sufficient condition:
$C_*(u)$ has a non-zero
limit along some sequence
$u_n\to-\infty$.

The saddle phase
$e^{-i\frac{2}{3}(|u|/3c_3)^{3/2}}$
oscillates with
irrationally-spaced
periods;
the Fresnel contribution
$C_0 e^{i\pi/4}\ne 0$
is constant.
If the saddle term is
smaller in amplitude
than the Fresnel term
(which requires a
quantitative bound on
the saddle amplitude
relative to the singularity),
then $\limsup|C_*(u)|
\ge|C_0|e^{i\pi/4}
-\sup|C_*^{(P)}|>0$.

\textbf{This is heuristic,
not a proof:}
the S+P split is not
rigorous as noted
in~\S\ref{sec:mechanisms}.
\end{remark}

\begin{problem}[Rigorous lower
  bound]
\label{prob:lower}
Prove Conjecture~\ref{conj:sharp}.
First step:
make the S+P partition rigorous
(cf.\ Rem.~\ref{rem:scale})
and compute the error
in the partition
uniformly in $u$.
Second step:
quantify the saddle
amplitude $|C_*^{(P)}(u)|$
and compare to $|C_0|$.
\end{problem}

\begin{remark}[Consequence
  for Paper~XV \Open]
\label{rem:consequence}
If Conjecture~\ref{conj:sharp}
holds in this model,
then Assumption~ass:Phi-cubic
of Paper~XV is false,
and $S(c)=\Theta(c^{\theta_0})$
(Ces\`aro is sharp).
This would be a
\emph{rigorous negative result}
for the PSWF compression
programme:
no improvement over
Ces\`aro is possible
via oscillatory methods
in this regime.
\end{remark}

\end{document}
